\documentclass{UoYCSproject}
\addbibresource{dummyBib.bib}
\author{Maciek Zaweracz}
\title{A Comparative Study on Error Correction Algorithms for LEO-Based CubeSats}
\date{20 April 2026}
\supervisor{Adrian Bors}
\BSc

%\dedication{To all students everywhere}

\acknowledgements{
I would like to take the opportunity to express my deepest gratitude to my 
supervisor, Dr Adrian Bors, for his invaluable feedback and support throughout this project.

I would also like to extend my sincere thanks to Dr Rupesh Kumar, whose expertise 
in quantum key distribution and insightful discussions greatly enhanced the depth and direction of this dissertation.

Additionally, I am deeply grateful to Killian Murphy for his guidance and practical support, 
particularly in the areas of system implementation and technical problem-solving, which were instrumental to the success of this work.

Finally, I would like to thank my family and friends for their love, support, and encouragement.
}

% More definitions & declarations in example.ldf

\begin{document}
\pagenumbering{roman}
\maketitle
\listoffigures
\listoftables
%\renewcommand*{\lstlistlistingname}{List of Listings}
%\lstlistoflistings

\begin{summary}

Secure digital communication depends on two parties --- conventionally
called Alice and Bob --- agreeing on a shared secret key without
leaking any of it to an adversary (Eve). Today this is done using
classical cryptography whose security rests on mathematical hardness
assumptions that sufficiently powerful quantum computers are expected
to break. Quantum Key Distribution (QKD) replaces those assumptions
with security grounded in quantum mechanics: any attempt by Eve to
measure the quantum states Alice sends to Bob inevitably disturbs
them, and the legitimate parties can detect this disturbance
\cite{grosshans_continuous_2002, laudenbach_continuous-variable_2018}.
Because optical fibre loses signal rapidly over long distances, global
QKD is expected to rely on satellites in Low Earth Orbit (LEO).
CubeSats --- small standardised spacecraft on the order of 10\,kg
\cite{oi_cubesat_2017} --- offer a low-cost entry point but impose
tight constraints on size, weight, power, onboard compute, memory and
classical communication bandwidth.

This dissertation is set in the context of the Satellite Platform for
Optical Quantum Communications (SPOQC), the Quantum Communications
Hub's 12U CubeSat that will demonstrate Continuous-Variable QKD
(CV-QKD) from space to a ground station at Heriot-Watt University
\cite{noauthor_hubs_nodate, simmons_hogs_2025}. CV-QKD is the variant
of QKD that uses standard telecom optics and therefore fits best into
a constrained satellite platform, but it also places the largest
burden on classical post-processing. The dominant post-processing cost
is \textit{information reconciliation}: the step in which Alice and
Bob, having ended up with noisy near-copies of the same bit-string,
correct their disagreements over a public classical channel without
handing their key to Eve \cite{mueller_performance_2025}. Two
reconciliation families dominate the literature. \textit{Low-Density
Parity-Check (LDPC)} codes are a one-way coding approach in which Bob
sends a compact parity message and Alice corrects locally.
\textit{Cascade} is an interactive protocol that locates errors
through repeated rounds of parity comparison and binary search over
the classical channel. Prior work characterises each algorithm in
isolation but rarely compares them under the combined pressure of
CV-QKD, very low Signal-to-Noise Ratio (SNR) and CubeSat-grade
hardware. This dissertation asks: under realistic LEO CubeSat
constraints, which of the two offers the better
performance-feasibility trade-off for CV-QKD?

The approach is in two parts. Part~I measures reconciliation behaviour
for three quasi-cyclic LDPC codes and the original 4-pass Cascade
protocol under matched pre-correction Bit Error Rate (BER), covering
the full range from the low-BER regime in which Cascade is
conventionally characterised through to the high-BER regime expected
on a free-space LEO link. A statistics-tracking contribution to the
Quantum Communications Hub's internal LDPC codebase emits
machine-readable output aligned with the open-source \texttt{Cascade-cpp}
reference, so that both algorithms feed the same analysis pipeline.
Part~II embeds those measurements in a first-order LEO mission model
at a representative low-SNR operating point
($\mathrm{SNR}\!\approx\!0.07$), scaled over three source rates (2, 10
and 100\,MHz), two candidate classical links (Radio-Frequency X-band
and optical), and three onboard compute paces.

Two structural findings emerge. First, LDPC and Cascade cross over in
reconciliation efficiency partway through the BER range ($q\!\approx\!0.22$): Cascade is
close to the Shannon limit at low BER, LDPC pulls ahead at higher BER,
and the LEO-relevant low-SNR regime sits firmly on LDPC's side.
Second, at the mission-relevant operating point Cascade leaks roughly
two orders of magnitude more classical information per reconciled key
bit than LDPC ($\sim$118 vs.\ $\approx$1) and does so across $\sim$20
sequential round-trips per frame against LDPC's single one-way
message. Plugged into the mission model, LDPC passes every bandwidth
and latency check but is bounded by onboard compute, which can be
closed either by deferring processing across orbits or by investigating
hardware acceleration of its inner loop. Cascade only remains feasible
under idealised per-frame parallelism on a low-latency optical
classical channel at modest source rate.

The research question is therefore answered conditionally: LDPC is the
preferred choice for a realistic LEO CubeSat CV-QKD mission under
radio-frequency or mixed classical links at the expected low SNR;
Cascade remains a plausible fallback only with a low-latency optical
classical channel at modest source rate and higher SNR. The practical
implication is that reconciliation strategy and classical-link choice
should be co-designed. Natural next steps are combining the measured
leakage with privacy amplification to estimate the end-to-end
secret-key rate, benchmarking both algorithms on a space-grade
processor such as the Q8 Xiphos planned for SPOQC
\cite{rao_quantum_2025}, and extending the comparison to
dynamic pass-level SNR and alternative protocol variants. No
significant legal, social, ethical, professional or commercial issues
were raised by this study.

\end{summary}

\begin{ethics}
This project raises no significant ethical, legal, social, professional or
commercial issues. It involves no human participants, no personal data, and
no classified material. 

Work with the internal SPOQC \texttt{cvqkd-reconciliation}
library was carried out under the supervision of the Quantum Communications
Hub team; the statistics-tracking contribution is intended to be released 
alongside the mission codebase. All other tools and
data used in this dissertation are open source, and the analysis pipeline,
configurations, and processed results produced here are reproducible from
the accompanying repository.
\end{ethics}

\chapter{Introduction}
\label{cha:Introduction}

\section{Context and Motivation}
\label{sec:intro-context}

Secure communication requires that two parties, Alice and Bob, agree on a
shared secret key without disclosing it to any adversary (Eve)
\cite{mackay_information_nodate}. Quantum Key Distribution (QKD) provides
such an agreement mechanism with security grounded in the laws of quantum
mechanics: any attempt by Eve to measure the transmitted states introduces
detectable disturbance \cite{grosshans_continuous_2002}. This dissertation
focuses on Continuous-Variable QKD (CV-QKD), which uses standard
telecom-compatible coherent optics and is therefore well suited to
size-, weight- and power-constrained satellite platforms; discrete-variable
alternatives are set out in Section~\ref{sec:bg-qkd}
\cite{laudenbach_continuous-variable_2018}.

Fibre-based QKD is loss-limited beyond a few hundred kilometres without
trusted relays or quantum repeaters \cite{lauterbach_study_2023,
neumann_quantum_2017, podmore_qkd_2021}, so satellite links --- particularly
in Low Earth Orbit (LEO) --- are a strategically important route towards
global quantum-secure networking, with CubeSats offering an especially
low-cost entry point \cite{oi_cubesat_2017, lauterbach_study_2023,
alvaro_space_2025, bouzoukis_overview_2025}. Representative
quantum-communications CubeSat concepts aim to consolidate payload and bus
hardware within an envelope on the order of 10\,kg total mass and a few
litres of volume \cite{oi_cubesat_2017}, at the cost of strict
size-, weight- and power (SWaP) budgets, limited onboard compute, finite
memory and time-varying link conditions \cite{oi_cubesat_2017,
sivasankaran_cubesat_2022}. These constraints interact directly with the
information reconciliation (IR) stage of QKD, in which Alice and Bob must
correct mismatches between their correlated raw keys over an authenticated
classical channel \cite{laudenbach_continuous-variable_2018,
mueller_performance_2025}. The specific mission context for this dissertation
is the Satellite Platform for Optical Quantum Communications (SPOQC), the
Quantum Communications Hub's 12U CubeSat scheduled to demonstrate in-orbit
QKD from space to the Hub Optical Ground Station (HOGS)
at Heriot-Watt University, Edinburgh, carrying a dual quantum payload with
matching ground receivers at HOGS to demonstrate two complementary QKD
approaches \cite{noauthor_hubs_nodate, simmons_hogs_2025}. The dissertation
is therefore motivated by a practical systems question: even when a
reconciliation algorithm is theoretically sound, is it operationally suitable
for a CV-QKD CubeSat mission under realistic LEO constraints?

\section{Problem and Research Question}
\label{sec:research-question}

Information reconciliation in CV-QKD can become the dominant practical
bottleneck in the low Signal-to-Noise Ratio (SNR) or high-noise regions
characteristic of free-space LEO links \cite{milicevic_quasi-cyclic_2018,
gumus_novel_2021, mueller_performance_2025}. The literature contains
substantial work on each candidate reconciliation algorithm in isolation ---
interactive Cascade and one-way Low-Density Parity-Check (LDPC) codes --- but
comparatively little direct, application-specific evidence for
reverse-reconciliation CV-QKD at very low SNR on a CubeSat-grade platform. 
This dissertation therefore addresses the following research question:

\begin{quote}
\textit{Under realistic LEO CubeSat constraints, which error-correction
algorithm - one-way LDPC-based reconciliation or interactive Cascade-style
reconciliation - provides the better performance-feasibility trade-off
for CV-QKD?}
\end{quote}

\noindent The study is limited to the reconciliation stage of
reverse-reconciliation CV-QKD in a LEO CubeSat context; a full spacecraft
design, detailed optical link model, and end-to-end cryptographic integration
are out of scope. The contribution is a focused review of reconciliation
methods for constrained space-based CV-QKD, a matched-condition simulation
comparison of LDPC and Cascade, a statistics-tracking extension to the
internal SPOQC \texttt{cvqkd-reconciliation} library that aligns its output
with \texttt{Cascade-cpp}, and a first-order feasibility model linking the
measurements to LEO CubeSat mission constraints.

\section{Structure of the Dissertation}
\label{sec:structure}

The dissertation is structured as follows:
\begin{itemize}
    \item Chapter~\ref{cha:background} introduces the prerequisite knowledge,
    reviews the literature and identifies the research gap;
    \item Chapter~\ref{cha:methodology} defines the evaluation framework in
    two parts --- Part~I, algorithmic comparison, and
    Part~II, first-order LEO CubeSat feasibility model --- and a final
    Implementation and Verification section;
    \item Chapter~\ref{cha:results} presents and analyses the results,
    mirroring the Part~I / Part~II structure with discussion and limitations; and
    \item Chapter~\ref{cha:conclusion} answers the research question with
    practical implications and directions for future work.
\end{itemize}

% =========================================================
\chapter{Background and Literature Review}
\label{cha:background}
% =========================================================

% ---------------------------------------------------------
\section{CubeSat Communication Systems}
\label{sec:bg-cubesat}
% ---------------------------------------------------------

CubeSats are standardised small satellites built from multiples of a
$10 \times 10 \times 10\,\text{cm}$ unit (1U) \cite{oi_cubesat_2017}. They
substantially reduce development time and launch cost compared with conventional
spacecraft, at the price of strict size, weight and power (SWaP) budgets that
limit payload volume, antenna size, thermal control, onboard memory and
processor capability \cite{oi_cubesat_2017, sivasankaran_cubesat_2022,
bouzoukis_overview_2025}. LEO operation adds a time-critical dimension: contact
with a fixed ground station occurs in short, intermittent passes whose link
quality varies with elevation, atmospheric path length, antenna geometry and
pointing accuracy \cite{oi_cubesat_2017, jordan-parra_evaluation_2024}. Reference
mission reports of university-led CubeSats in the 2U--6U class
commonly observe around three passes per day of at most 12 minutes each
above a minimum-elevation threshold, giving a narrow operational
envelope for any onboard processing task
\cite{nieto-peroy_cubesat_2019}. Space
QKD missions require both a quantum optical link and a classical communication
link, so any reconciliation method must fit within practical limits on pass
duration, bandwidth, latency, onboard compute and memory --- constraints that
are particularly binding on 6U--12U platforms \cite{simmons_hogs_2025}.

% % ---------------------------------------------------------
% \section{Information Theory}
% \label{sec:bg-info-theory}
% % ---------------------------------------------------------

% Information theory, established by Claude Shannon, formalises how uncertainty is quantified,
% how noise affects communication, and how much information must be exchanged
% to recover a message reliably \cite{mackay_information_nodate, shannon1948mathematical}. For a binary
% variable with error probability $p$, the binary entropy function
% \begin{equation}
% H_2(p) = -p \log_2 p - (1-p)\log_2(1-p)
% \end{equation}
% gives the asymptotic lower bound on the number of bits that must be revealed
% to correct errors on a binary symmetric channel (BSC) \cite{shannon1948mathematical}. In QKD every such bit
% is assumed known to the eavesdropper and therefore reduces the final secret-key
% length. In binary QKD the post-sifting channel is typically modelled as a BSC;
% in CV-QKD the underlying physical channel is continuous-valued and more
% naturally described by a Gaussian noise model
% \cite{laudenbach_continuous-variable_2018}. This distinction matters for
% reconciliation: CV-QKD carries soft reliability information that one-way codes
% such as LDPC can exploit, whereas interactive protocols such as Cascade operate
% on hard binary decisions \cite{leverrier_multidimensional_2008,
% milicevic_quasi-cyclic_2018, martinez-mateo_demystifying_2015,
% reis_quantum_2019}.

% ---------------------------------------------------------
\section{Quantum Key Distribution}
\label{sec:bg-qkd}
% ---------------------------------------------------------

QKD establishes a shared secret key between Alice and Bob using quantum states
as information carriers. Its security rests on the fact that any eavesdropper
measurement introduces detectable disturbance
\cite{grosshans_continuous_2002, laudenbach_continuous-variable_2018}.

The two main QKD families are discrete-variable (DV) and continuous-variable
(CV). DV-QKD encodes information in discrete quantum states (e.g.\ photon
polarisation or phase) and typically relies on single-photon detectors, which
are difficult to accommodate on constrained satellite platforms. CV-QKD instead
encodes information in the continuous quadrature amplitudes of coherent states
and recovers them using homodyne or heterodyne detection with standard telecom
optics, avoiding the DV detector infrastructure
\cite{laudenbach_continuous-variable_2018}. This dissertation focuses on
Gaussian-Modulated Coherent-State (GMCS) CV-QKD: Alice modulates coherent states
according to a Gaussian distribution and Bob measures them by homodyne
detection, which avoids squeezed or entangled light and supports high-speed,
multi-bit-per-pulse operation
\cite{grosshans_continuous_2002, grosshans_quantum_2003}. Practical GMCS implementations have demonstrated
operation at source rates of hundreds of MHz using commercially available
GHz-bandwidth balanced homodyne detectors
\cite{tang_performance_2020}, bracketing the 2--100\,MHz source-rate range
considered in this work.

\subsection{The CV-QKD Process and Reverse Reconciliation}
\label{subsec:bg-cvqkd-process}

A CV-QKD pipeline consists of six stages: Alice transmits Gaussian-modulated
coherent states; Bob performs coherent detection; a subset of the data is used
for parameter estimation; the continuous observations are mapped to binary
variables using multidimensional reconciliation
\cite{leverrier_multidimensional_2008}; an information reconciliation protocol
removes mismatches between Alice's and Bob's correlated keys; and privacy
amplification compresses the reconciled key. Unlike DV-QKD, CV-QKD does not
produce a binary raw key directly, which makes reconciliation particularly
demanding. In \emph{direct} reconciliation Bob corrects his data to match
Alice; in \emph{reverse} reconciliation Alice corrects her data to match Bob.
Reverse reconciliation is necessary in high-loss channels such as space-to-ground
links because it preserves a positive key rate where direct reconciliation
fails \cite{grosshans_continuous_2002, grosshans_quantum_2003}. It fixes the
direction of the classical communication and places the decoding task onboard
the satellite --- the intersection at which reconciliation cost meets CubeSat
mission constraints \cite{mueller_performance_2025}.

% =========================================================
% \chapter{Literature Review}
% \label{cha:literature}
% =========================================================

\section{Space QKD and Past CubeSat Missions}
\label{sec:lit-space}

Satellite QKD is motivated by the distance limitations of terrestrial optical
fibre. 
While fibre-based QKD is effective over metropolitan and inter-city
scales, global secure communication requires methods that avoid the high
attenuation of long fibre links. Free-space optical channels therefore provide
an important route towards long-distance quantum-secure communication
\cite{lauterbach_study_2023}.

A major milestone in the space QKD area was the \textit{Micius} satellite, which
demonstrated long-distance satellite-to-ground QKD, at a distance of up to 1200km, and established the
practical viability of space-based quantum links
\cite{lauterbach_study_2023, lu_micius_2022}. It was the world's first quantum satellite mission, 
and although it differed significantly from a regular CubeSat platform, weighing over 600kg, 
it showed that quantum communication can be sustained over free-space paths far beyond the range 
typically achievable in fibre.

CubeSat-oriented work instead considers whether space QKD can be realised
within the mass, volume, power, and operational limits of a nanosatellite.
The CQuCoM concept makes these constraints explicit through a 6U CubeSat in
an approximately 10\,kg, 6\,L envelope, with platform power around 11\,W, a
30\,Wh battery, a 90\,mm transmit telescope, and quantum passes limited to
night-time eclipse opportunities of about 6 minutes on average
\cite{oi_cubesat_2017}. This shows that feasibility depends not only on the
optical protocol, but also on whether pointing, power storage, and pass
duration can be made compatible with quantum transmission. Canada's Quantum Encryption and Science Satellite (QEYSSat), 
set to launch in 2026, aims to address the same problem from a different systems perspective \cite{agency_quantum_2017}. 
It uses a LEO trusted-node architecture with the satellite primarily as a receiver for
ground-to-space uplinks, shifting source complexity and much of the processing
burden to the ground segment \cite{podmore_qkd_2021}.

These studies establish the feasibility and strategic relevance of space
QKD, but they do not directly answer how a resource-constrained CubeSat
should select a reconciliation strategy when compute, memory, bandwidth and
latency are treated as mission-level constraints
\cite{jordan-parra_evaluation_2024}. This issue is especially important for
SPOQC, which provides a concrete CubeSat-oriented CV-QKD mission context. Because CV-QKD requires continuous
measurement outcomes to be converted into shared binary secret keys, the
choice of reconciliation protocol becomes a central systems question. Recent
SPOQC/HOGS material indicates losses of up to around 50~dB and key rates on
the order of 100~bits/s, suggesting that mission feasibility may depend not
only on optical performance but also on the cost and behaviour of the
reconciliation stage \cite{noauthor_hubs_nodate, simmons_hogs_2025}.

\section{Information Reconciliation Approaches}
\label{sec:lit-reconciliation}

QKD reconciliation is dominated by two approaches: interactive protocols such
as Cascade, and code-based protocols such as LDPC, each with distinct
trade-offs in efficiency, communication structure and compute cost
\cite{martinez-mateo_demystifying_2015}.

Cascade is conceptually simple: errors are located through repeated block
partitioning, parity comparison and binary search
\cite{martinez-mateo_demystifying_2015, reis2019quantum}, and its cost and
leakage depend strongly on channel conditions and configuration
\cite{reis2019quantum, tupkary_using_2023}. Mart\'inez-Mateo et al.\ report
$f_{\mathrm{EC}} = 1.09013$ at a pre-correction Bit Error Rate (BER) of
$Q = 2\%$ ($\mathrm{SNR}\!\approx\!6.0$) and $f_{\mathrm{EC}} = 1.09541$ at
$Q = 5\%$ ($\mathrm{SNR}\!\approx\!3.5$) for $10^{4}$-bit frames, at the
price of tens to over a hundred channel uses per frame depending on the
Quantum BER (QBER) \cite{martinez-mateo_demystifying_2015}. At higher BER,
improved efficiency typically costs substantially more interaction
\cite{martinez-mateo_demystifying_2015}, which matters here because low-SNR
CV-QKD can produce effective binary channels in regimes where communication
overhead and leakage become major practical costs. Refinements such as
Pacher et al.'s feedback protocol with power-of-two parity weights
reduce leakage to $\eta_{\mathrm{IR}}\!<\!1.03$ across the full BSC range
\cite{pacher_information_2015}, but only the original 4-pass variant is
evaluated here (see Section~\ref{sec:limitations}).

LDPC reconciliation replaces repeated two-way parity exchange with one-way
syndrome transmission and local iterative decoding. It can approach the
Shannon limit while reducing classical communication overhead, and in CV-QKD
is especially prominent because it can exploit soft channel information,
making it well suited to multidimensional reconciliation and low-SNR settings
\cite{leverrier_multidimensional_2008, milicevic_quasi-cyclic_2018}. Slice
reconciliation with LDPC has been used to push CV-QKD efficiency above 90\%
at SNR as low as 0.03 \cite{jouguet_high-bit-rate_2014}, confirming the
relevance of the low-SNR regime targeted in this work. The literature
contains many variants, including blind and rate-adaptive methods
for uncertain channels \cite{gumus_novel_2021, mueller_performance_2025},
and other one-way code families such as rateless spinal codes have been
proposed as alternatives reaching $\sim\!96.5\%$ efficiency at similarly low
SNR \cite{wen_novel_2019}; these all share the same basic profile of high
efficiency and low communication overhead, but with performance that depends
on code selection, large block lengths and sufficient compute
\cite{milicevic_quasi-cyclic_2018, gumus_novel_2021}.

\section{Comparative Literature and Research Gap}
\label{sec:lit-comparative}

Direct comparisons between Cascade and LDPC are less common than studies of
each protocol individually, and existing ones depend strongly on channel
model, QKD protocol, hardware assumptions and chosen metrics. Mueller et
al.\ show that communication rounds, verification cost, classical-channel
behaviour and robustness to parameter mismatch can all change the relative
attractiveness of a protocol, and that protocol assessment should be
application-oriented --- their own comparison is restricted to DV-QKD and
$Q \le 0.10$ ($\mathrm{SNR}\!\approx\!2.0$), which does not cover the low-SNR
CV-QKD regime relevant here \cite{mueller_performance_2025,
jordan-parra_evaluation_2024}. A second consistent lesson is that successful
reconciliation does not automatically imply useful key generation: a protocol
may correct errors while revealing too much information, or requiring too much
communication, for a secret key to remain after privacy amplification
\cite{mueller_performance_2025, tupkary_using_2023}.

Four aspects of this literature define the gap addressed by this dissertation.
First, most comparative reconciliation work considers DV-QKD or simplified
binary channel models rather than reverse-reconciliation CV-QKD with
soft-information preprocessing. Second, the low-SNR regime characteristic of
free-space CV-QKD is not where Cascade is usually most attractive. Third,
existing studies rarely translate reconciliation behaviour into a CubeSat
mission model with explicit limits on pass duration, bandwidth, latency,
memory and onboard compute. Fourth, many of the strongest LDPC results assume
hardware acceleration or very large block lengths that are not obviously
representative of a CPU-only CubeSat platform
\cite{milicevic_quasi-cyclic_2018, mueller_performance_2025,
jordan-parra_evaluation_2024, simmons_hogs_2025}. The methodology in
Chapter~\ref{cha:methodology} is designed to close this gap.

% =========================================================
\chapter{Methodology}
\label{cha:methodology}
% =========================================================

This dissertation uses a quantitative, simulation-based methodology to compare
LDPC and Cascade under controlled channel conditions and then assess their
feasibility in a representative LEO CubeSat mission context. The methodology
has two parts. Part~I defines the theoretical comparison framework --- channel
model, protocol pipelines, and metrics --- under matched error conditions.
Part~II maps those metrics to first-order communication, compute, and memory
constraints relevant to an onboard implementation. A final
Section~\ref{sec:meth-implementation} describes the software artefacts,
experiment orchestration and verification procedures used to produce the
measurements.

Simulation is appropriate because the target mission environment is not
available as a controlled testbed, the relevant operating region spans a range
of low-SNR conditions, and the two algorithms cannot be compared fairly without
systematic control of channel difficulty
\cite{mueller_performance_2025, jordan-parra_evaluation_2024}.

% ---------------------------------------------------------
\section{Part I: Simulated Performance Analysis}
\label{sec:meth-part1}
% ---------------------------------------------------------

LDPC and Cascade operate on different internal representations. LDPC is a
one-way, soft-decision code-based scheme used here within a CV-QKD pipeline.
Cascade is an interactive, hard-decision protocol evaluated on a Binary
Symmetric Channel (BSC), a standard abstraction in
which each transmitted bit is independently flipped with some fixed
probability. They therefore cannot be compared on strictly identical
internal terms.

The comparison basis used in this dissertation is the measured
\textit{pre-correction BER}. On the LDPC side, SNR is swept through the full
CV-QKD simulation path and the resulting hard decisions are used to measure
the actual BER seen before decoding. On the Cascade side, the requested BSC
error rate is adjusted so that the measured BER aligns with the LDPC
operating points. This preserves the CV-QKD interpretation on the LDPC side while
allowing cross-algorithm comparison under matched effective channel difficulty
\cite{mueller_performance_2025}.

\subsection{CV-QKD Channel Model}
\label{subsec:meth-channel}

The LDPC experiments are performed within a simplified CV-QKD simulation path
based on an Additive White Gaussian Noise (AWGN) channel:
\begin{equation}
Y = \sqrt{t}\,X + N, \qquad N \sim \mathcal{N}(0,\sigma_n^2),
\end{equation}
where $X$ is Alice's transmitted quadrature, $Y$ is Bob's received quadrature,
$t$ is effective channel transmittance, and $N$ is Gaussian noise. This AWGN
form is the standard first-order abstraction for a free-space optical link, as it
compresses physical loss into a single multiplicative factor $\sqrt{t}$ and all
noise sources into a single Gaussian term, which is sufficient to drive
reconciliation at a chosen operating point. The corresponding operating SNR is
\begin{equation}
\text{SNR} = \frac{t \sigma_X^2}{\sigma_n^2},
\end{equation}
i.e.\ the ratio of received signal power to noise power, the canonical way
of expressing how difficult a link is for the decoder: the lower the SNR,
the noisier the samples that reconciliation has to turn into an agreed
bit-string.

This is a fixed-SNR benchmark model rather than a full pass-level optical
simulation. In a real LEO pass, SNR varies with elevation, pointing, and
background conditions. For reconciliation benchmarking, however, a fixed
operating point is an appropriate first-order abstraction because it isolates
algorithmic behaviour from link-budget dynamics
\cite{milicevic_quasi-cyclic_2018, jordan-parra_evaluation_2024}.

\subsection{LDPC and Cascade Pipelines}
\label{subsec:meth-pipelines}

For LDPC, each experiment follows the same sequence:
\begin{enumerate}
    \item Gaussian sample generation;
    \item additive noise injection at the selected SNR;
    \item 8-dimensional multidimensional reconciliation preprocessing;
    \item hard-decision bit mapping and syndrome generation;
    \item Min-Sum LDPC decoding.
\end{enumerate}

The multidimensional reconciliation stage follows the Leverrier--Grangier 8D
procedure, which maps Gaussian-valued observations to binary variables and
produces both hard decisions and reliability information for decoding
\cite{leverrier_multidimensional_2008}.

Cascade is evaluated separately on a binary symmetric channel. Each run applies
the standard cycle of shuffling, block partitioning, parity exchange, binary
search, and look-back correction \cite{martinez-mateo_demystifying_2015}. Cascade is therefore not driven directly by
SNR, but by the measured BER used to align it to LDPC operating points.

\subsection{Recorded Metrics and Statistical Treatment}
\label{subsec:meth-metrics}

Seven quantities are recorded. Reconciliation efficiency measures how close
the algorithm operates to the Shannon limit and is defined as $\eta = \mathrm{reconciliation\_bits}/(n \cdot H_2(q))$; here
$H_2(q) = -q\log_2 q - (1-q)\log_2(1-q)$ is the binary entropy function,
which gives the minimum number of bits per symbol
that must be disclosed to reconcile a BSC with error probability $q$, so
$\eta = 1$ is the Shannon limit and larger values indicate more leakage. 
Pre-correction BER ($q$) serves both as a sanity check and as
the alignment variable for joint LDPC/Cascade comparison. Remaining BER and
FER quantify post-correction reliability; FER matters for LDPC because
frame-level decoding failure is still possible.
Reconciliation bits per key bit normalises classical communication overhead
across frame sizes. Elapsed process and real time per key bit feed the Part~II
compute model. Cascade message counts capture sequential parity exchanges for
latency modelling. LDPC decoder iterations proxy decoding difficulty. Means
and sample standard deviations are computed over $10^{4}$ independent trials
per operating point; timing and communication metrics are normalised per key
bit to compare 256--2048-bit LDPC frames fairly against 10\,000-bit Cascade
keys \cite{martinez-mateo_demystifying_2015}.

% ---------------------------------------------------------
\section{Part II: System-Level Feasibility Model}
\label{sec:meth-part2}
% ---------------------------------------------------------

\subsection{Mission Abstraction and Operating Point}
\label{subsec:meth-system-model}

Part~II converts the Part~I per-frame measurements into first-order resource
requirements under representative LEO CubeSat constraints. It uses an algebraic
check at one mission-relevant operating point; residual BER,
FER, efficiency and post-privacy-amplification key rate are out of scope
(See section~\ref{sec:future-work}).

\begin{figure}[!htbp]
    \centering
    \includegraphics[width=0.78\linewidth]{Overleaf-Dissertation-Files/Figures/overpass_fig.pdf}
    \caption{Overpass geometry used in the Part~II mission abstraction.
    The CubeSat traces an arc above the Optical Ground Station (OGS); the
    QKD-feasible window corresponds to a bounded elevation range symmetric
    about zenith.}
    \label{fig:overpass}
\end{figure}

% The nominal SPOQC QKD window is taken here between
%     $\theta = -30^{\circ}$ (QKD start) and $\theta = +30^{\circ}$ (QKD end),
%     during which the quantum signal is transmitted and Alice buffers the
%     corresponding classical data onboard for later reconciliation.

The model assumes a Sun-synchronous LEO CubeSat at 500--600\,km with one
high-quality optical pass per day over a fixed ground station, yielding a
120\,s quantum window (Figure~\ref{fig:overpass}) in which Alice (satellite)
buffers her transmitted data and Bob (ground) measures and computes the side
information; reconciliation runs in a separate classical window. Three source-rate cases are considered:
2\,MHz (the SPOQC rate \cite{rao_quantum_2025}), 10\,MHz and 100\,MHz (upper-bound stress). 
Two classical links are considered: an optical classical sidechannel
$T_{\text{opt}} \approx 30$\,s at $\sim$1\,Gbps, and a Radio-Frequency (RF)
X-band window $T_{\text{RF}} \approx 300$\,s at $\sim$50\,Mbps
\cite{noauthor_hubs_nodate}. Onboard memory is capped at
$M_{\max} \approx 125$\,GB. Free-space CV-QKD on a LEO link is expected to operate in the low-SNR
region of roughly $[0.006, 0.07]$
\cite{milicevic_quasi-cyclic_2018, noauthor_hubs_nodate}, which through
the empirical 8D mapping (Section~\ref{subsec:meth-verification})
corresponds to pre-correction BER near $q \approx 0.40$. All numerical
claims are evaluated at $q = 0.40$ ($\mathrm{SNR}\!\approx\!0.07$), which
sits at the upper edge of this window and therefore represents a
conservatively high-BER (worst-case-leakage) operating point; $q = 0.10$
($\mathrm{SNR}\!\approx\!2.0$) is retained only as a higher-SNR
sensitivity reference.

\subsection{Time and Communication Model}
\label{subsec:meth-time-model}

Let $f_{\text{src}}$ be the source rate and $T_q = 120$\,s the quantum
window. Per-pass workload is
\begin{equation}
N = f_{\text{src}} \cdot T_q,
\end{equation}
i.e.\ the number of raw quantum symbols that must eventually be reconciled
from a single quantum pass; it sets the scale of every other resource demand
below. Completion time is modelled to first order as
\begin{equation}
T_{\text{total}} \approx T_{\text{compute}} + T_{\text{tx}} + T_{\text{latency}},
\end{equation}
splitting end-to-end cost into three additive contributions: onboard
processing, classical-link transmission, and the waiting time accumulated
over sequential round-trips. Required classical volume is
\begin{equation}
B = N \cdot r,
\end{equation}
where $r$ is the reconciliation bits per key bit measured in Part~I; $B$ is
therefore the total public-channel traffic a reconciliation run generates,
and is what must fit into the available link budget.
Bandwidth feasibility is $B \le R \cdot T_{\text{window}}$ for the
relevant link rate $R$ and window
$T_{\text{window}} \in \{T_{\text{opt}},T_{\text{RF}}\}$, expressing the
simple requirement that the protocol cannot send more bits than the link
can carry in the time available.
Cascade's sequential interaction latency is
\begin{equation}
T_{\text{latency}} = N_{\text{seq}}\,\tau,
\end{equation}
the product of the number of sequential parity round-trips per frame
$N_{\text{seq}}$ (measured in Part~I) and the per-round-trip delay $\tau$;
this is the latency penalty interactive protocols pay on top of raw
transmission time. LDPC has
$N_{\text{seq}} = 1$ by construction. Whether $T_{\text{latency}}$ is
paid once per frame (fully concurrent Cascade engines, \emph{pipelined}
bound) or across all $N/n$ frames (\emph{serial} bound) is treated as
a bounding question in
Section~\ref{subsec:results-latency-feasibility}.

\subsection{Latency Scenarios}
\label{subsec:meth-latency}

Cascade is simulated without link latency; $\tau$ is applied analytically.
Four cases are considered: ideal baseline: $\tau = 0$\,s, optimistic RF:
$\tau = 0.1$\,s, conservative RF: $\tau = 0.5$\,s and optical classical:
$\tau \approx 5$\,ms (the light-speed Round-Trip Time, RTT, at LEO
altitude \cite{conde_round_2024}).

\subsection{Compute Scenarios}
\label{subsec:meth-compute-scenarios}

Part~I timings were measured on an Apple M1 Pro. Both reference
implementations are scalar and single-threaded per reconciliation: the
LDPC binary uses a single C process, and the Cascade core algorithm is
single-threaded C++17. Neither uses the GPU,
Neural Engine, explicit SIMD or multi-threading within one decode, so
the figure reflects best-case single-core CPU behaviour rather than the
peak capability of the silicon. Rescaling these timings to a
specific flight processor without dedicated benchmarking would be
speculative, so three generic compute-pace scenarios are used instead:
\begin{itemize}
    \item \textbf{Accelerated} ($0.1\times$ dev time) --- assisted by a
    Field-Programmable Gate Array (FPGA) or Digital Signal Processor (DSP),
    which are hardware platforms well suited to LDPC's regular inner loop.
    \item \textbf{Nominal} ($1\times$) --- directly reuse the Part~I timings.
    \item \textbf{Slow} ($10\times$) --- a round-number stand-in for a
    lower-clock, in-order, CPU-only space-grade Single-Board Computer
    (SBC) such as the Q8 Xiphos used on SPOQC; treated as a one-decade
    envelope rather than a Q8-specific benchmark.
\end{itemize}
Compute is checked against two windows: the classical window after the
pass, and an extended multi-orbit budget
$T_{\text{orbit}} \approx 24$\,h representing deferred processing across
orbits. The former requires compute to finish before the key is
downlinked; the latter accepts that compute runs to the next
mission-relevant trigger, at the cost of a delayed key.

\subsection{Feasibility Criteria}
\label{subsec:meth-feasibility}

Each (algorithm, source-rate) combination is judged against four
criteria:
\begin{itemize}
    \item \textbf{Bandwidth}: $B \le R \cdot T_{\text{window}}$ for the relevant link;
    \item \textbf{Latency}: $T_{\text{latency}} \le T_{\text{window}}$ for each $\tau$ scenario;
    \item \textbf{Compute}: $T_{\text{compute}} \le T_{\text{window}}$ for each pace scenario;
    \item \textbf{Memory}: $N \cdot b_{\text{sym}} \le M_{\max}$, with $b_{\text{sym}} \approx 4$\,B per buffered sample.
\end{itemize}
These produce the pass/fail summary in
Section~\ref{subsec:results-feasibility-matrix}.

% ---------------------------------------------------------
\section{Implementation and Verification}
\label{sec:meth-implementation}
% ---------------------------------------------------------

The framework above is turned into measurements by a short software stack
built around two libraries, a structured-output layer and an external
experiment runner, all checked against published references before use.

\subsection{Software Implementations}
\label{subsec:meth-implementations}

LDPC is implemented using the \texttt{cvqkd-reconciliation} C library
developed for the SPOQC CV-QKD payload
\cite{noauthor_hubs_nodate, noauthor_uoy-researchcvqkd-reconciliation_nodate}. Using a
mission-relevant codebase makes the measurements more representative of the
intended deployment context than a standalone implementation. The decoder
uses Min-Sum decoding on quasi-cyclic parity-check matrices. A matrix named
\texttt{MxN\_zK} denotes a base parity-check matrix of shape $M \times N$
expanded by a circulant lifting factor $z = K$, producing a final
parity-check matrix of size $(Mz) \times (Nz)$ and therefore a frame length
of $n = Mz$ bits per reconciliation. Three matrices are evaluated:
\begin{itemize}
    \item \texttt{256x255\_z1} — frame length $n = 256$ bits,
    \item \texttt{512x511\_z4} — frame length $n = 2048$ bits,
    \item \texttt{1024x1023\_z1} — frame length $n = 1024$ bits.
\end{itemize}
These are existing low-SNR matrices available in the SPOQC codebase and were
selected as the study set for comparative testing. In the remainder of the
dissertation LDPC variants are identified by their effective frame length
(256, 1024, and 2048 bits) rather than by base-matrix name, since the frame
length is what controls per-frame behaviour, statistics, and normalised
metrics.

Cascade is implemented using the open-source \texttt{Cascade-cpp} library
\cite{rijsman_brunorijsmancascade-cpp_2026}. This implementation records
timing, message-count, parity-request, and leakage statistics, which makes
it suitable both for direct comparison and for later latency modelling.

\subsection{Statistics Tracking Contribution}
\label{subsec:meth-tracking}

A key implementation contribution of this dissertation is the addition of a
dedicated statistics-tracking system to the internal SPOQC LDPC repository.
The original \texttt{cvqkd-reconciliation} library was extended so that each
operating point emits structured statistics in Newline-Delimited JSON
(NDJSON) format. The output schema was designed to align closely with
\texttt{Cascade-cpp}, so that both algorithms could be parsed and analysed
through the same downstream workflow. The added tracking records the
quantities defined in Section~\ref{subsec:meth-metrics}. For LDPC, these
values are accumulated online over repeated independent frames using
Welford's algorithm, which provides numerically stable running means and
sample standard deviations without retaining all raw samples in memory. 
This is useful for large sweeps and high
trial counts, and it keeps the measurement layer lightweight enough to
remain compatible with the existing SPOQC implementation. The tracked
output fields for both implementations are summarised in
Appendix~\ref{chap:appendix-output}.

\subsection{Experiment Orchestration}
\label{subsec:meth-orchestration}

LDPC sweeps are orchestrated using an external Python runner driven by a
YAML configuration file. The runner varies SNR and parity-check matrix
selection, invokes the LDPC binary, and stores machine-readable outputs for
later analysis. A geometric SNR schedule is used in the range
$\mathrm{SNR} \in [10^{-4}, 7]$ so that the low-SNR regime is sampled more
densely, since this is the expected operating range for a LEO CubeSat and
the behaviour changes most rapidly there \cite{milicevic_quasi-cyclic_2018}.
Through the empirical 8D mapping, this SNR window corresponds to measured
pre-correction BER in approximately $[0.015, 0.499]$; LDPC data is therefore
absent from plots below $q \approx 0.015$. A representative configuration
is provided in Appendix~\ref{chap:appendix-config}. Cascade operating points
are generated separately over requested BSC error rates and later matched
to LDPC points using the measured pre-correction BER.

\subsection{Verification}
\label{subsec:meth-verification}

Each implementation is checked against known references before being used in
the comparative study. Cascade is verified by reproducing the efficiency
behaviour of the published \texttt{Cascade-cpp} reference and the
Mart\'inez-Mateo et al.\ literature baseline
\cite{martinez-mateo_demystifying_2015, mueller_performance_2025}. LDPC is
verified by reproducing the internal SPOQC $1 - 2\cdot\mathrm{BER}$ reference
curve, which exercises the Gaussian generator, multidimensional reconciliation
preprocessor, syndrome generator and Min-Sum decoder together.
A separate alignment check compares the measured pre-correction BER values
used in the matched LDPC/Cascade plots; any systematic mismatch here would
indicate an error in the comparison procedure rather than in the
reconciliation methods themselves. A final alignment check compares the
measured pre-correction BER against the Binary Phase-Shift Keying (BPSK)
reference curve $\tfrac{1}{2}\mathrm{erfc}(\sqrt{\mathrm{SNR}})$, the
canonical SNR-to-BER relationship for a hard-decision binary channel with
Gaussian noise; the difference between the two is due to the 8-dimensional
(8D) reconciliation preprocessor used in our experiments,
so we use the measured mapping throughout our results. Supporting figures
are in Appendix~\ref{chap:appendix-verification}.

% =========================================================
\chapter{Results and Analysis}
\label{cha:results}
% =========================================================

\section{Part I: Algorithmic Results}
\label{sec:results-part1}

Part~I runs LDPC and Cascade through the matched-BER protocol of
Section~\ref{sec:meth-part1} and reports what each algorithm actually
does under the same channel noise level. Five questions are then answered:
\emph{how close to the information-theoretic leakage limit does
each run} (efficiency), \emph{how often does the output key contain
residual errors} (reliability), \emph{how much classical traffic does
each generate} (overhead), \emph{how long does each take on a CPU}
(compute), and \emph{how interactive is each} (messages per frame).
Together these feed directly into the Part~II feasibility model. 

For readability, comparison plots use pre-correction BER as the primary
x-axis and a secondary SNR axis on top. The SNR axis uses the empirical
LDPC mapping measured in the CV-QKD pipeline (see
Section~\ref{subsec:meth-verification} and
Appendix~\ref{chap:appendix-verification}). Each point is a
mean over $10^4$ independent trials with ±1$\sigma$ shaded bands where
shown.

\FloatBarrier
\subsection{Reconciliation Efficiency}
\label{subsec:results-efficiency}

Efficiency $\eta$ (defined in Section~\ref{subsec:meth-metrics}) measures
how much more information each algorithm leaks over the classical channel
than the information-theoretic minimum; $\eta = 1$ is the Shannon floor
and smaller is better. It matters because this leakage is what the
secret-key rate ultimately has to absorb.
Figures~\ref{fig:efficiency-full}, \ref{fig:efficiency-zoom-2} and
\ref{fig:efficiency-zoom} show reconciliation efficiency
$\eta = \mathrm{reconciliation\_bits} / (n \cdot H_2(q))$ against matched
pre-correction BER; the top SNR axis is the empirical LDPC mapping measured
in the CV-QKD pipeline. Three observations follow from these curves.

First, in the low-BER regime canonically used to evaluate Cascade
\cite{martinez-mateo_demystifying_2015}, Cascade is close to the Shannon limit
while LDPC is heavily over-provisioned. Cascade reaches
$\eta \approx 1.14$ at $q = 0.01$, $1.19$ at $q = 0.05$
($\mathrm{SNR}\!\approx\!3.5$), and $1.21$ at $q = 0.10$
($\mathrm{SNR}\!\approx\!2.0$); LDPC, whose syndrome is fixed by code rate,
produces $\eta \approx 6.8\text{--}8.1$ at $q = 0.02$
($\mathrm{SNR}\!\approx\!6.0$). The Cascade values lie a few percent above
Mart\'inez-Mateo et al.'s 1.09 / 1.10 at $Q = 2\% / 5\%$
\cite{martinez-mateo_demystifying_2015}, consistent with using the original
4-pass variant without BICONF and Cascade-cpp's default block schedule \cite{martinez-mateo_demystifying_2015}.

Second, the two curves cross at $q \approx 0.22$
($\mathrm{SNR}\!\approx\!0.66$). Beyond the crossover LDPC is the more
efficient: $\eta = 1.14, 1.03, 1.00$ at $q = 0.30, 0.40, 0.49$
($\mathrm{SNR}\!\approx\!0.30, 0.07, 0.001$), against $\eta = 1.40, 1.50, 1.59$
for Cascade at the same points. The expected low-SNR LEO region sits past
the crossover, so LDPC has the leakage advantage there.

Third, the LDPC curve asymptotes to $0.996\text{--}0.999$ as $q \to 0.5$.
This is not super-Shannon: at $q = 0.5$, $H_2(q) \to 1$ and
$\eta \to (M-1)/M \approx 0.996$ by construction. The apparent sub-unit
efficiency coincides with decoder failure
(Section~\ref{subsec:results-reliability}), so the efficiency metric has
simply become uninformative there.

The three frame lengths produce visually indistinguishable efficiency curves
except in the extreme low-BER tail, where the fractional syndrome-size
differences $(M-1)/M$ matter. The step discontinuities in the Cascade curve
near $q = 0.09, 0.14, 0.18, 0.25, 0.37$ correspond to Cascade-cpp's piecewise
block-size schedule switching to a larger $k_1$
\cite{martinez-mateo_demystifying_2015, rijsman_brunorijsmancascade-cpp_2026}.

\begin{figure}[!htbp]
    \centering
    \includegraphics[width=0.80\linewidth]{experiments/graphs/results/efficiency_full.png}
    \caption{Reconciliation efficiency $\eta$ over the full BER range.
    Crossover at $q \approx 0.22$ ($\mathrm{SNR}\!\approx\!0.66$); LDPC
    asymptotes to $\eta \approx 1$ but this coincides with decoder failure
    (cf.\ Figure~\ref{fig:remaining-fer}).}
    \label{fig:efficiency-full}
\end{figure}

\begin{figure}[!htbp]
    \centering
    \begin{subfigure}[t]{0.49\linewidth}
        \centering
        \includegraphics[width=\linewidth]{experiments/graphs/results/efficiency_zoom_2.png}
        \caption{High-SNR regime.}
        \label{fig:efficiency-zoom-2}
    \end{subfigure}\hfill
    \begin{subfigure}[t]{0.49\linewidth}
        \centering
        \includegraphics[width=\linewidth]{experiments/graphs/results/efficiency_zoom.png}
        \caption{Low-SNR regime.}
        \label{fig:efficiency-zoom}
    \end{subfigure}
    \caption{Zoomed views of reconciliation efficiency $\eta$ in two operating regimes. Figure~\subref{fig:efficiency-zoom-2} 
    shows the high-SNR region ($q \le 0.10$, $\mathrm{SNR}\!\gtrsim\!2.0$), where Cascade is conventionally characterised; 
    the LDPC curve is omitted because it remains off-scale ($\eta > 1.3$ throughout). Figure~\subref{fig:efficiency-zoom} 
    shows the low-SNR region more relevant to the expected free-space LEO CV-QKD operating range; 
    here Cascade is omitted because it remains off-scale ($\eta \approx 1.5\text{--}1.6$ across the window).}
    \label{fig:efficiency-zooms}
\end{figure}

\FloatBarrier
\subsection{Correction Reliability}
\label{subsec:results-reliability}

Correction reliability asks how often each algorithm hands back a key
that is actually correct --- residual BER for bit-level errors
and FER for whole-frame failures. A protocol that leaks little but
frequently returns the wrong key is useless, so this sets the limit 
of where each algorithm can be operated.
Figures~\ref{fig:remaining-ber} and~\ref{fig:remaining-fer} show
post-correction BER and FER.

Cascade exhibits a well-defined low-BER failure mode: at $q = 0.0002$
(two bit errors in a $10^4$-bit frame) FER reaches 46\%, falling to 4.5\% at
$q = 0.001$ and below $2\times10^{-4}$ by $q = 0.01$. Residual BER in the
failed frames is tiny ($\sim\!10^{-5}$) because each undetected failure
leaves only one or two bits flipped. This is the two-error ambiguity of the
original 4-pass variant without BICONF: two errors in blocks of matching
parity cannot be localised by binary search, so Cascade terminates on a
parity-consistent but incorrect key \cite{martinez-mateo_demystifying_2015}.
Above $q \approx 0.01$, FER is $\le 10^{-4}$ throughout the measured range.

LDPC shows the complementary waterfall: residual BER and FER are both zero
for $q \le 0.38$ ($\mathrm{SNR}\!\gtrsim\!0.10$) and rise sharply above it.
At $q = 0.45$ ($\mathrm{SNR}\!\approx\!0.017$) FER is 3.6\% ($n\!=\!1024$),
7.4\% ($n\!=\!256$) and 8.0\% ($n\!=\!2048$); at $q = 0.49$
($\mathrm{SNR}\!\approx\!0.001$) FER is 38\%, 38\% and 89\% for
$n = 256, 1024, 2048$. Two points deserve emphasis. \emph{(i)} The longest
frame $n\!=\!2048$ fails most often at extreme BER, which is counter-intuitive
if one expects block length to improve thresholds; the explanation is that
this code is a $4\times$ lift of the 512-column base matrix
\texttt{512x511\_z4}, and for QC-LDPC in CV-QKD the near-threshold behaviour
is dominated by the base-matrix structure rather than the lifted length
\cite{milicevic_quasi-cyclic_2018, mueller_performance_2025}. \emph{(ii)} The
wide $\pm 1\sigma$ bands past $q = 0.4$ reflect bimodal behaviour: frames
either decode fully or fail catastrophically (residual BER $\approx 0.2$), as
expected of Min-Sum past the decoder threshold
\cite{mackay_information_nodate}.

\begin{figure}[!htbp]
    \centering
    \includegraphics[width=0.95\linewidth]{experiments/graphs/results/remaining_ber.png}
    \caption{Post-correction BER. Cascade is $\approx 0$ above $q = 0.01$;
    LDPC shows a clean waterfall past $q \approx 0.38$
    ($\mathrm{SNR}\!\approx\!0.10$).}
    \label{fig:remaining-ber}
\end{figure}

\begin{figure}[!htbp]
    \centering
    \includegraphics[width=0.95\linewidth]{experiments/graphs/results/remaining_fer.png}
    \caption{Post-correction frame error rate. The narrow Cascade spike below
    $q = 0.01$ is the two-error ambiguity failure of original-variant
    Cascade \cite{martinez-mateo_demystifying_2015}; the LDPC waterfall past
    $q = 0.38$ is decoder-threshold behaviour
    \cite{milicevic_quasi-cyclic_2018}.}
    \label{fig:remaining-fer}
\end{figure}

\FloatBarrier
\subsection{Communication Overhead}
\label{subsec:results-overhead}

Communication overhead is the number of classical bits exchanged for
each reconciled key bit; this is what has to fit through the satellite's
classical link budget in Part~II, so it is the single metric that most
directly governs bandwidth feasibility.
Figure~\ref{fig:recon-bits} shows reconciliation bits per key bit, the
metric that most sharply separates the two algorithms. LDPC is flat at
$(M-1)/M$, giving 0.996, 0.999 and 0.9980 for $n = 256, 1024, 2048$,
independent of BER because the syndrome size is fixed by the parity-check
matrix. Cascade grows from 7.7 at $q = 0.01$ to 46.1 at $q = 0.10$
($\mathrm{SNR}\!\approx\!2.0$), 75.5 at $q = 0.20$
($\mathrm{SNR}\!\approx\!0.79$), 117.8 at $q = 0.40$
($\mathrm{SNR}\!\approx\!0.07$) and 128.4 at $q = 0.49$
($\mathrm{SNR}\!\approx\!0.001$). At the high-BER end Cascade therefore
transmits $\sim\!130\times$ more classical information per reconciled key bit
than LDPC; even in the traditional Cascade regime ($q < 0.1$) the ratio is
8--47$\times$ in LDPC's favour. This is the qualitative behaviour expected
of any iterative parity-exchange protocol
\cite{martinez-mateo_demystifying_2015}, but the magnitude is the driver of
the bandwidth-feasibility analysis in Part~II.

\begin{figure}[!htbp]
    \centering
    \includegraphics[width=0.95\linewidth]{experiments/graphs/results/recon_bits_per_key_bit.png}
    \caption{Classical communication overhead per reconciled key bit. LDPC is
    effectively constant at $\approx 1$ bit/key bit; Cascade grows to $\sim\!130$
    bits/key bit at $q = 0.5$.}
    \label{fig:recon-bits}
\end{figure}

\FloatBarrier
\subsection{Compute Cost}
\label{subsec:results-compute}

Compute cost is the processing time per reconciled key bit and controls
whether reconciliation keeps up with the quantum source on limited
onboard hardware. It feeds the three compute-pace scenarios of
Section~\ref{subsec:meth-compute-scenarios} and determines whether
compute is the binding constraint at the mission-relevant operating
point.
Figures~\ref{fig:real-time} and~\ref{fig:process-time} show wall-clock
and CPU time per reconciled key bit on the development laptop (Apple M1
Pro);
the two agree within 1--2\% at every operating point, so no significant I/O
or scheduling overhead is being captured. Both algorithms lie in
$1\text{--}8 \times 10^{-7}\,\mathrm{s/bit}$. Cascade is fastest per bit
throughout, from $3\times10^{-8}\,\mathrm{s/bit}$ at $q = 0.01$ to
$3.6\times10^{-7}$ at $q = 0.49$ ($\mathrm{SNR}\!\approx\!0.001$); LDPC sits
at $2.0\text{--}6.5\times10^{-7}\,\mathrm{s/bit}$, with the $n\!=\!2048$
variant cheapest at low BER but climbing sharply above $q = 0.4$
($\mathrm{SNR}\!\approx\!0.07$) as mean Min-Sum iterations rise from
$\approx 3$ to $\approx 12$ (Figure~\ref{fig:iterations}). These numbers
should not be read as ``Cascade is cheaper onboard'': they compare two
different codebases on a general-purpose CPU, and LDPC's sparse
belief-propagation inner loop parallelises naturally on DSPs and FPGAs
whereas Cascade's serial binary search and sequential parity exchange do
not. Absolute per-bit times are used in Part~II only through the three
compute-pace scenarios of Section~\ref{subsec:meth-compute-scenarios}.

\begin{figure}[!htbp]
    \centering
    \includegraphics[width=0.95\linewidth]{experiments/graphs/results/real_time_per_bit.png}
    \caption{Wall-clock time per reconciled key bit on a desktop-class CPU.}
    \label{fig:real-time}
\end{figure}

\begin{figure}[!htbp]
    \centering
    \includegraphics[width=0.95\linewidth]{experiments/graphs/results/process_time_per_bit.png}
    \caption{Process (CPU) time per reconciled key bit. Tracks
    Figure~\ref{fig:real-time} to within 1--2\%, confirming no significant
    I/O or scheduling overhead in the measurement.}
    \label{fig:process-time}
\end{figure}

\begin{figure}[!htbp]
    \centering
    \includegraphics[width=0.95\linewidth]{experiments/graphs/results/iterations.png}
    \caption{Mean Min-Sum iterations per LDPC frame, with the Cascade
    reference value of 4 passes shown for comparison. The $n\!=\!2048$
    variant's iteration count rises sharply past $q \approx 0.4$
    ($\mathrm{SNR}\!\approx\!0.07$), driving the compute-cost increase in
    Figure~\ref{fig:real-time}.}
    \label{fig:iterations}
\end{figure}

\FloatBarrier
\subsection{Interaction Cost}
\label{subsec:results-interaction}

Interaction cost is the number of sequential round-trips a frame
requires. LDPC needs one by construction, Cascade many, and each
round-trip costs one $\tau$ of link latency in Part~II --- so this is
the metric that decides whether Cascade stays within the classical
window under a given link.
Figure~\ref{fig:Cascade-messages} shows the number of sequential
parity-request round-trips per Cascade reconciliation. The mean ranges from
15 to 45 with a saw-tooth shape: each discontinuity is a Cascade-cpp
block-schedule transition that raises the initial block size $k_1$ and
temporarily reduces request count \cite{martinez-mateo_demystifying_2015,
rijsman_brunorijsmancascade-cpp_2026}. The count is $\approx 45$ at
$q = 0.01$, $\approx 37$ at $q = 0.10$ ($\mathrm{SNR}\!\approx\!2.0$), and
never falls below $\approx 15$ even at $q = 0.5$. LDPC transmits a single
classical block per frame independently of BER, so the 15--45$\times$ gap in
sequential message count drives Part~II's latency analysis: whatever the
round-trip delay $\tau$, Cascade pays it at least 15 times per 10\,000-bit
frame and LDPC pays it once.

\begin{figure}[!htbp]
    \centering
    \includegraphics[width=0.95\linewidth]{experiments/graphs/results/Cascade_messages.png}
    \caption{Cascade sequential parity-request messages per reconciliation.
    Saw-tooth discontinuities mark block-schedule transitions in Cascade-cpp
    \cite{rijsman_brunorijsmancascade-cpp_2026}. LDPC is structurally one
    message per frame and is not shown.}
    \label{fig:Cascade-messages}
\end{figure}

\FloatBarrier
\subsection{Synthesis of Part I Results}
\label{subsec:results-synthesis}

Table~\ref{tab:results-synthesis} summarises the five Part~I metrics at three
representative operating points covering the measured range.

\begin{table}[H]
    \centering
    \small
    \setlength{\tabcolsep}{5pt}
    \renewcommand{\arraystretch}{1.15}
    \begin{tabularx}{\textwidth}{
        >{\raggedright\arraybackslash}p{0.24\textwidth}
        >{\centering\arraybackslash}X
        >{\centering\arraybackslash}X
        >{\centering\arraybackslash}X
    }
        \toprule
        Metric & $q = 0.02$ (SNR$\approx$6.0) & $q = 0.20$ (SNR$\approx$0.79) & $q = 0.45$ (SNR$\approx$0.017) \\
        \midrule
        Efficiency $\eta$
            & 6.94 / 1.16
            & 1.38 / 1.29
            & 1.01 / 1.54 \\

        Residual FER
            & 0 / $<10^{-3}$
            & 0 / 0
            & 0.036 / 0 \\

        Recon.\ bits / key bit
            & 0.999 / 13.5
            & 0.999 / 75.5
            & 0.999 / 123.7 \\

        Messages / frame
            & --- / 45.4
            & --- / 27.8
            & --- / 19.7 \\

        Wall-clock time per key bit
        ($10^{-7}$\,s)
            & 2.2 / 0.5
            & 3.8 / 2.1
            & 5.4 / 3.4 \\
        \bottomrule
    \end{tabularx}
    \caption{Representative Part~I values at low, moderate and high BER. Entries are given as LDPC / Cascade.}
    \label{tab:results-synthesis}
\end{table}

Three quantitative conclusions follow, each an input to Part~II:
\begin{enumerate}
    \item \textbf{Efficiency crossover at $q \approx 0.22$
        ($\mathrm{SNR}\!\approx\!0.66$).} LDPC holds the leakage advantage in
        the LEO-relevant regime; the ``Cascade is close to Shannon'' narrative
        applies only to the low-BER end.
    \item \textbf{Classical-volume ratio $10^{1}\text{--}10^{2}$ in LDPC's
        favour.} Even at $q = 0.10$ ($\mathrm{SNR}\!\approx\!2.0$), Cascade
        discloses $\approx 47$ bits per reconciled bit against LDPC's
        $\approx 1$.
    \item \textbf{Sequential-message ratio $15\text{--}45\times$ in LDPC's
        favour.} BER-insensitive to within a factor of three; dominates
        feasibility whenever $\tau$ is non-trivial.
\end{enumerate}
Working against LDPC: non-zero residual BER/FER past $q \approx 0.38$
($\mathrm{SNR}\!\approx\!0.10$) with bimodal failure that discards a
double-digit fraction of frames at worst-case BER, and $\approx\!2\text{--}3\times$
higher per-bit CPU time than Cascade on the test machine at low-to-moderate
BER. Working against Cascade: a non-zero undetected-error rate in the
very-low-BER regime, intrinsic to the original 4-pass variant
\cite{martinez-mateo_demystifying_2015}. Neither failure mode is fatal, but
both must be carried into the system-level model.

\section{Part II: Feasibility Analysis}
\label{sec:results-part2}

Part~II applies the mission model of Section~\ref{sec:meth-part2} to the
Part~I measurements at one representative operating point $q = 0.40$
($\mathrm{SNR}\!\approx\!0.07$) and at source rates
$f_{\text{src}} \in \{2, 10, 100\}$\,MHz. The higher-SNR reference
$q = 0.10$ ($\mathrm{SNR}\!\approx\!2.0$) is used for sensitivity only. All
inputs come directly from Part~I; no additional simulation is performed.

\subsection{Workload}
\label{subsec:results-workload}

Workload sets the scale every other feasibility check is measured
against: how many raw samples are produced per pass and how much RAM
they occupy onboard. Per-pass workload $N = f_{\text{src}} \cdot T_q$ is
the number of raw quantum samples buffered for reconciliation; buffered
data is $D = N \cdot b_{\text{sym}}$ with $b_{\text{sym}} = 4$\,B per
sample, assuming each is stored as a 32-bit float. For downstream
normalisation, each sample is treated as one reconciled bit after 8D
preprocessing, so $N$ is used interchangeably for samples and key bits
at this level of abstraction. Results are in
Table~\ref{tab:part2-workload}. Memory demand stays below the 125\,GB
budget at every source rate, although at 100\,MHz the 48\,GB buffer
already consumes roughly 40\% of the budget before the syndrome,
rotation metadata and double-buffering overhead of a real pipeline are
counted.

\begin{table}[H]
    \centering
    \small
    \begin{tabular}{lrrrr}
        \toprule
        $f_{\text{src}}$ & $N$ (samples/pass) & $D$ (GB) & LDPC ($n{=}1024$) & Cascade ($n{=}10^4$)\\
        \midrule
        2\,MHz   & $2.4\!\times\!10^{8}$  & 0.96 & $2.3\!\times\!10^{5}$ & $2.4\!\times\!10^{4}$\\
        10\,MHz  & $1.2\!\times\!10^{9}$  & 4.8  & $1.2\!\times\!10^{6}$ & $1.2\!\times\!10^{5}$\\
        100\,MHz & $1.2\!\times\!10^{10}$ & 48   & $1.2\!\times\!10^{7}$ & $1.2\!\times\!10^{6}$\\
        \bottomrule
    \end{tabular}
    \caption{Per-pass workload for $T_q = 120$\,s and $b_{\text{sym}} = 4$\,B.
    $N$ is raw samples (treated as key bits post-8D). LDPC and Cascade columns show the number of frames needed per sample per pass. }
    \label{tab:part2-workload}
\end{table}

\subsection{Bandwidth Feasibility}
\label{subsec:results-comm-feasibility}

Bandwidth feasibility checks whether the classical traffic a protocol
generates fits through the available link in the time the link is open.
Required classical volume $B = N \cdot r$ is evaluated with
$r_{\text{LDPC}} \approx 1.00$ and
$r_{\text{Cascade}} \approx 118$ bits/key bit at $q = 0.40$
($\mathrm{SNR}\!\approx\!0.07$; Section~\ref{subsec:results-overhead}).
Capacities are $R_{\text{opt}} \cdot T_{\text{opt}} \approx 30$\,Gbit and
$R_{\text{RF}} \cdot T_{\text{RF}} \approx 15$\,Gbit.

\begin{table}[H]
    \centering
    \small
    \begin{tabular}{lrrr}
        \toprule
        Algorithm (and $r$) & 2\,MHz & 10\,MHz & 100\,MHz\\
        \midrule
        LDPC ($r \approx 1$)                  & 0.24\,Gbit & 1.2\,Gbit  & 12\,Gbit  \\
        Cascade at $q{=}0.40$ ($r \approx 118$) & 28\,Gbit   & 141\,Gbit  & 1410\,Gbit\\
        Cascade at $q{=}0.10$ ($r \approx 46$)  & 11\,Gbit   & 55\,Gbit   & 550\,Gbit \\
        \bottomrule
    \end{tabular}
    \caption{Required classical reconciliation volume $B$. Capacities:
    RF $\approx 15$\,Gbit; optical $\approx 30$\,Gbit.}
    \label{tab:part2-bandwidth}
\end{table}

LDPC fits comfortably on both links at every source rate. At $q = 0.40$
($\mathrm{SNR}\!\approx\!0.07$), Cascade exceeds the RF window at every rate
and the optical window from 10\,MHz onwards. At the higher-SNR reference
$q = 0.10$ ($\mathrm{SNR}\!\approx\!2.0$), only the 2\,MHz + optical
combination is bandwidth-feasible. This follows directly from the Part~I
finding that Cascade's per-bit overhead grows with BER by one to two orders
of magnitude (Section~\ref{subsec:results-overhead}).

\subsection{Latency Feasibility}
\label{subsec:results-latency-feasibility}

Latency feasibility asks whether Cascade's sequential round-trips can
complete inside the classical window; this is where the
``many-messages-per-frame'' result from Part~I turns into pass/fail on
a given link. LDPC is insensitive to $\tau$ so only Cascade is stressed
here.
Cascade's per-frame sequential message count at $q = 0.40$
($\mathrm{SNR}\!\approx\!0.07$) is $N_{\text{seq}} \approx 20$
(Section~\ref{subsec:results-interaction}). The pipelined bound treats all
frames as concurrent, giving $T_{\text{lat}} = N_{\text{seq}}\,\tau$; the
serial bound treats them frame-by-frame, giving
$(N/n) \cdot N_{\text{seq}}\,\tau$.

\begin{table}[H]
    \centering
    \small
    \begin{tabular}{lrrr}
        \toprule
        Channel & $\tau$ & Pipelined $T_{\text{lat}}$ & Serial $T_{\text{lat}}$ at 10\,MHz\\
        \midrule
        Optical              & 5\,ms    & 0.1\,s  & $\approx 12\,000$\,s ($\sim$\,3.3\,h)\\
        RF, optimistic       & 100\,ms  & 2\,s    & $\approx 2.4\!\times\!10^{5}$\,s ($\sim$\,2.8\,d)\\
        RF, conservative     & 500\,ms  & 10\,s   & $\approx 1.2\!\times\!10^{6}$\,s ($\sim$\,14\,d)\\
        \bottomrule
    \end{tabular}
    \caption{Cascade interaction latency at $q = 0.40$
    ($\mathrm{SNR}\!\approx\!0.07$), $N_{\text{seq}} \approx 20$. LDPC has
    $N_{\text{seq}} = 1$ by construction and is therefore insensitive to $\tau$.}
    \label{tab:part2-latency}
\end{table}

Under the pipelined bound Cascade is feasible in every case. Under the
serial bound it is feasible only for $\tau \to 0$, i.e.\ effectively the
optical case. Realistic RF links fail the serial bound by orders of
magnitude, so Cascade's practicality on a CubeSat depends on how much
frame-level parallelism can be sustained on flight hardware.

\subsection{Compute Feasibility}
\label{subsec:results-compute-feasibility}

Compute feasibility asks whether the onboard processor can finish
reconciling a full pass before the key is needed downstream. It is the
constraint where LDPC's regular, parallelisable inner loop is most
likely to bind on a CPU-only platform.
Using the Part~I per-bit timings at $q = 0.40$
($\mathrm{SNR}\!\approx\!0.07$): $\approx 5\!\times\!10^{-7}$\,s/bit (LDPC)
and $\approx 3\!\times\!10^{-7}$\,s/bit (Cascade), with
$T_{\text{compute}} = N \cdot t$.

\begin{table}[H]
    \centering
    \small
    \begin{tabular}{llrrr}
        \toprule
        $f_{\text{src}}$ & Algorithm & Accelerated ($0.1\times$) & Nominal ($1\times$) & Slow ($10\times$)\\
        \midrule
        2\,MHz   & LDPC    & 12\,s   & 120\,s   & 1200\,s   \\
                 & Cascade & 7\,s    & 72\,s    & 720\,s    \\
        10\,MHz  & LDPC    & 60\,s   & 600\,s   & 6000\,s   \\
                 & Cascade & 36\,s   & 360\,s   & 3600\,s   \\
        100\,MHz & LDPC    & 600\,s  & 6000\,s  & 60\,000\,s \\
                 & Cascade & 360\,s  & 3600\,s  & 36\,000\,s \\
        \bottomrule
    \end{tabular}
    \caption{Total onboard compute time at $q = 0.40$
    ($\mathrm{SNR}\!\approx\!0.07$) for the three compute-pace scenarios.}
    \label{tab:part2-compute}
\end{table}

At nominal pace LDPC already exceeds $T_{\text{RF}} = 300$\,s at 10\,MHz and
needs multi-orbit deferral at 100\,MHz. The slow scenario pushes both
algorithms beyond the post-pass window at 10\,MHz and beyond the 24\,h
inter-orbit budget at 100\,MHz. The accelerated scenario restores LDPC to
within the RF window up to 10\,MHz, the regime in which FPGA- or DSP-assisted
belief-propagation is worth investigating. Compute is therefore the dominant
constraint for LDPC whenever a key must be delivered within a single
post-pass window; the natural mitigations are multi-orbit deferral (accepting
a delayed key) or hardware acceleration of the LDPC inner loop
(Section~\ref{subsec:results-compute}).

\subsection{Feasibility Summary}
\label{subsec:results-feasibility-matrix}

The summary matrix draws the four per-criterion checks above together
into a single pass/fail table, making it easy to read off which
(algorithm, source-rate, link) combinations remain viable end-to-end.
Table~\ref{tab:part2-matrix} combines the bandwidth, latency and compute
checks at $q = 0.40$ ($\mathrm{SNR}\!\approx\!0.07$). Memory is omitted
because it does not bind, and BER/FER/efficiency are out of scope (Part~I,
Section~\ref{sec:future-work}).

\begin{table}[H]
    \centering
    \small
    \setlength{\tabcolsep}{2.5pt}
    \renewcommand{\arraystretch}{1.1}
    \begin{tabular}{>{\raggedright\arraybackslash}p{0.34\linewidth}|ccc|ccc}
        \toprule
        & \multicolumn{3}{c|}{LDPC} & \multicolumn{3}{c}{Cascade}\\
        Criterion & 2\,MHz & 10\,MHz & 100\,MHz & 2\,MHz & 10\,MHz & 100\,MHz\\
        \midrule
        Bandwidth, RF (15\,Gbit)          & \checkmark & \checkmark & \checkmark & --        & --        & --\\
        Bandwidth, optical (30\,Gbit)     & \checkmark & \checkmark & \checkmark & \checkmark & --        & --\\
        Latency, pipelined (any $\tau$)   & \checkmark & \checkmark & \checkmark & \checkmark & \checkmark & \checkmark\\
        Latency, serial RF ($\tau{=}0.5$\,s) & \checkmark & \checkmark & \checkmark & --        & --        & --\\
        Compute, nominal, 300\,s window   & \checkmark & --        & --        & \checkmark & --        & --\\
        Compute, nominal, 24\,h window    & \checkmark & \checkmark & \checkmark & \checkmark & \checkmark & \checkmark\\
        \bottomrule
    \end{tabular}
    \caption{Pass (\checkmark) / fail (--) feasibility summary at
    $q = 0.40$ ($\mathrm{SNR}\!\approx\!0.07$). RF capacity uses X-band
    50\,Mbps $\times$ 300\,s; optical uses 1\,Gbps $\times$ 30\,s. LDPC has
    $N_{\text{seq}} = 1$ and is therefore independent of both the
    pipelined/serial distinction and the choice of $\tau$.}
    \label{tab:part2-matrix}
\end{table}

Three observations follow from Table~\ref{tab:part2-matrix}:
\begin{enumerate}
    \item \textbf{LDPC is compute-bound.} All bandwidth and latency
    criteria pass; only compute pressures the short classical window.
    This can be relaxed either by multi-orbit deferral or, potentially,
    by hardware acceleration.
    \item \textbf{Cascade is communication- and latency-bound.}
    Bandwidth fails on RF at every source rate and on optical from
    10\,MHz upwards. Latency fails the serial bound on any realistic RF
    link.
    \item \textbf{Link choice matters more for Cascade than for LDPC.}
    Moving from RF to optical does not change LDPC's feasibility, but
    it is the difference between Cascade being infeasible and marginally
    feasible at low source rate.
\end{enumerate}

The \emph{kind} of failure matters more than the fact of it. At 10\,MHz
with an RF classical link at nominal compute pace, both algorithms
exceed the 300\,s post-pass window, but the dominant term of
$T_{\text{total}} = T_{\text{compute}} + T_{\text{tx}} +
T_{\text{latency}}$ (Section~\ref{subsec:meth-time-model}) is different.
LDPC's $\approx\!624$\,s is compute-dominated ($600$\,s compute,
$24$\,s to stream $1.2$\,Gbit over $50$\,Mbps, zero latency because
$N_{\text{seq}}=1$), while Cascade's $\approx\!3190$\,s is
transmission-dominated ($\sim\!47$\,min just to push $141$\,Gbit over
the same link). In mission-design terms, LDPC's overrun is a deferrable
backlog: one extra orbit clears it on the same hardware. Cascade's is a
link-class deficit: no amount of deferral on the RF link delivers a
key, so the choice is between upgrading the classical channel and
lowering the source rate.

LDPC is the stronger candidate under an RF classical link or at the
expected low SNR, while Cascade is only plausible alongside a
low-latency optical classical channel at modest source rate and higher
SNR. Passing the four criteria, however, does not imply a positive secret-key rate
after privacy amplification, and either algorithm --- Cascade in
particular --- may still leave zero usable
key once leakage is accounted for. A full secret-key-rate evaluation is
out of scope (limitations, Section~
\ref{sec:limitations}; future work, Section~\ref{sec:future-work}).

\section{Discussion}
\label{sec:discussion}

LDPC wins the expected LEO CubeSat regime (low SNR,
RF or mixed classical link); Cascade only wins in the narrow corner of
higher SNR on a low-latency optical classical channel at modest source
rate. The remainder of this section explains why.

Algorithm choice is governed by two structural facts: LDPC consumes
compute, and Cascade consumes classical communication and latency. Part~I
confirmed both in isolation; Part~II showed that under realistic LEO CubeSat
budgets they translate into opposite failure modes. Quantitatively, the two
efficiency curves crossed at $q \approx 0.22$
($\mathrm{SNR}\!\approx\!0.66$), past which LDPC is the more
leakage-efficient; at the mission-relevant operating point
$q = 0.40$ ($\mathrm{SNR}\!\approx\!0.07$), Cascade's classical volume
and sequential-message count exceeded LDPC's by one to two orders of
magnitude ($r \approx 118$ vs.\ $\approx 1$ bit/key bit, and
$N_{\text{seq}} \approx 20$ vs.\ $1$). Plugged into the mission model
these numbers translate into the four feasibility checks: LDPC passes
every bandwidth and latency check and is pressured only by compute,
which relaxes either through multi-orbit deferral or through
hardware acceleration of the belief-propagation inner loop; Cascade
fails RF bandwidth at every source rate considered and fails the
serial-latency bound on any realistic RF link, leaving only the
low-rate optical-classical corner feasible.

Because the expected low-SNR LEO region sits past the Part~I efficiency
crossover, the leakage conclusion from Part~I and the feasibility
conclusion from Part~II agree. The two practical implications for
mission design are that reconciliation strategy and classical-link
choice should be co-designed --- RF-only missions plan on LDPC,
optical-classical missions preserve both options --- and that onboard
compute, not bandwidth, is the realistic bottleneck for LDPC, so the
compute budget and the in-pass vs.\ cross-orbit completion policy
should be fixed at mission definition.

\section{Limitations}
\label{sec:limitations}

The study is first-order: fixed-SNR operating points rather than a dynamic
pass-level optical simulation, BER alignment rather than a shared internal
CV-QKD channel for Cascade, analytical latency rather than a network
emulator, and timings on development hardware rather than flight hardware.
It is also an error-correction feasibility study, not a security study:
leakage and the post-privacy-amplification secret-key rate are out of scope,
so a protocol that corrects every bit may still be cryptographically
unusable if it leaks too much. The LDPC conclusions --- including the
$q \approx 0.22$ ($\mathrm{SNR}\!\approx\!0.66$) crossover --- are tied to
the three fixed quasi-cyclic matrices tested and would shift for
rate-adaptive, blind, or alternatively lifted codes; Cascade is similarly
evaluated only in its original 4-pass variant.

Four further first-order assumptions in the Part~II feasibility model are
worth naming explicitly. \emph{(i) Symmetric RF link.} The RF bandwidth
check treats ``the RF link'' as a single 50\,Mbps X-band downlink, but
SPOQC's classical architecture is asymmetric: the X-band downlink is
$\sim$50\,Mbps, whereas the S-band uplink is around 76.8\,kbps, roughly
650$\times$ narrower \cite{noauthor_hubs_nodate}. Under reverse
reconciliation, Bob (ground) sends the LDPC syndrome and Cascade parity
requests upward, so the uplink is the genuine bottleneck for both
algorithms; a realistic traffic-directional analysis would tighten every
``RF-feasible'' entry in the feasibility matrix and is left as future work.
\emph{(ii) Compute platform.} Both Part~I timings were produced by
       scalar, single-threaded reference implementations running on one
       performance core of an Apple M1 Pro, with no GPU or Neural Engine; 
       this penalises LDPC more than Cascade because
       Min-Sum belief propagation is highly parallel whereas Cascade is
       inherently serial, so the reported LDPC/Cascade compute gap is a
       conservative upper bound on LDPC's real deployment cost. The
       $10\times$ ``slow'' scenario is a rough stand-in for a space-grade
       single-board computer rather than a measured benchmark on flight
       hardware. \emph{(iii) Time-of-pass dynamics.} The first-order
model defines $T_{\text{total}} = T_{\text{compute}} + T_{\text{tx}} +
T_{\text{latency}}$ and evaluates each term against a single post-pass
window, without scheduling multiple passes per day, computing how many
orbits of classical downlink are needed to deliver a key, or modelling
the gaps between passes. Where the compute or bandwidth term exceeds the
300\,s post-pass window, the model falls back to the 24\,h inter-orbit
budget as an upper bound rather than computing the exact number of
deferred orbits. \emph{(iv) Memory model.} The buffered-sample estimate
uses 4\,B per symbol (32-bit float) as a conservative upper bound on the
soft-sample representation, and does not include rotation metadata,
syndrome storage, or the double-buffering that a real pipeline would
require.

None of these caveats invalidate the qualitative conclusions --- LDPC
compute-bound, Cascade communication- and latency-bound --- but the
matrix entries are better read as ``closes to first order'' than as
precise mission budgets.

\chapter{Conclusions and Future Work}
\label{cha:conclusion}

\section{Summary and Implications}
\label{sec:summary}

This dissertation asked whether LDPC-based or Cascade reconciliation
is the better fit for a CV-QKD mission flown on a LEO CubeSat. It
answered the question in two steps. Part~I compared the two algorithms
on the same channel and measured how much classical information each
leaks, how reliable each is, how fast each runs, and how many round-trips
each needs. Part~II placed those measurements inside a simple model of
a real satellite pass --- finite time over a ground station, a finite
classical link, finite onboard compute and memory --- and asked which
combinations actually close.

Under the radio-frequency classical link
that a CubeSat of this class is most likely to fly with, and in the
low-SNR regime that a free-space link spends most of its time in,
LDPC is the stronger choice: its classical traffic and round-trip
count are small and fixed, so the only pressure left is onboard
compute, which can be absorbed either by processing across orbits or
by investigating hardware acceleration of the decoder. Cascade remains
an option, but only in a narrower corner of the design space: a
low-latency optical classical channel, a modest source rate, and a
higher-SNR operating point. Outside that corner its back-and-forth
interaction simply does not fit inside the time and bandwidth a pass
leaves for classical communication. A practical consequence for
mission design is that reconciliation strategy and classical-link
choice should be chosen together rather than independently.

Zooming out, the comparison is a small step in a much larger picture.
Classical cryptography currently relies on mathematical problems that
future quantum computers are expected to solve, so every long-lived
encrypted conversation today is, in principle, already at risk. QKD
replaces those assumptions with security grounded in physics, and
CubeSat-scale satellites are the most affordable path to making it work
at continental and global scales. The results here suggest that the
post-processing stage is one of the main levers that decides whether such
a mission closes at all. Getting the reconciliation-and-link pairing
right on a small, power-limited spacecraft is therefore a concrete
enabler of cheaper, more general, and more widely deployable
quantum-secure communication. The direct next steps are to close the
remaining loop to an end-to-end secret-key rate, and to repeat this
analysis on representative flight hardware so that the compute
bottleneck identified here can be quantified rather than bounded.

\section{Future Work}
\label{sec:future-work}

Three extensions follow directly from this scope. 

\textbf{End-to-end
secret-key rate}: combine Part~I leakage with privacy amplification and a
channel model to verify positivity of the final key. 

\textbf{Flight-hardware
benchmarking}: replace the generic compute-pace scenarios with measured
timings on a space-grade processor such as the Q8 Xiphos used on SPOQC, and
investigate whether FPGA or DSP offload of the LDPC inner loop can realise
the accelerated scenario in practice. 

\textbf{Dynamic channel model and
protocol variants}: replace fixed-SNR analysis with an elevation-dependent
pass simulation, and extend the comparison to rate-adaptive LDPC and
alternative Cascade variants --- for example BICONF (an extension that
confirms reconciliation with additional randomised parity checks) --- 
to see whether the regime in which Cascade is preferable
widens.

\appendix

\chapter{Verification Plots}
\label{chap:appendix-verification}

This appendix collects the verification and alignment plots supporting
Section~\ref{subsec:meth-verification}.

\begin{figure}[!htbp]
    \centering

    % Row 1: Cascade
    \begin{minipage}[t]{0.32\textwidth}
        \centering
        \includegraphics[width=\linewidth]{experiments/graphs/verification/cascade/VerificationCascadeOriginal0.png}
        \subcaption{Cascade efficiency reproduced in this dissertation workflow.}
    \end{minipage}\hfill
    \begin{minipage}[t]{0.32\textwidth}
        \centering
        \includegraphics[width=\linewidth]{experiments/graphs/verification/cascade/Cascadec_cpp_author_fig_1.png}
        \subcaption{\texttt{Cascade-cpp} reference efficiency graph.}
    \end{minipage}\hfill
    \begin{minipage}[t]{0.32\textwidth}
        \centering
        \includegraphics[width=\linewidth]{experiments/graphs/verification/cascade/fig1_demyst_ir_Cascade.png}
        \subcaption{Published efficiency behaviour from \textit{Demystifying the Information Reconciliation Protocol Cascade}.}
    \end{minipage}

    \vspace{1em}

    % Row 2: LDPC
    \begin{minipage}[t]{0.48\textwidth}
        \centering
        \includegraphics[width=\linewidth]{experiments/graphs/verification/ldpc/1-2xBER.png}
        \subcaption{Reproduced LDPC $1 - 2\cdot\mathrm{BER}$ curve.}
    \end{minipage}\hfill
    \begin{minipage}[t]{0.48\textwidth}
        \centering
        \includegraphics[width=\linewidth]{experiments/graphs/verification/ldpc/ed_ldpc_1-2xBER.png}
        \subcaption{Internal SPOQC reference $1 - 2\cdot\mathrm{BER}$ curve.}
    \end{minipage}

    \caption{Verification figures used to validate the Cascade and LDPC implementations. 
    Panels (a)--(c) compare the reproduced Cascade efficiency behaviour against software 
    and literature references; panels (d)--(e) compare the 
    reproduced LDPC preprocessing behaviour against the internal reference baseline.}
    \label{fig:verification_appendix}
\end{figure}

\begin{figure}[!htbp]
    \centering
    \includegraphics[width=0.78\linewidth]{experiments/graphs/verification/actual_bit_error_rate.png}
    \caption{Measured pre-correction BER alignment used in the matched-condition
    LDPC/Cascade comparison. Agreement supports the correctness of the
    alignment independently of the reconciliation algorithms themselves.}
    \label{fig:ber_alignment_appendix}
\end{figure}

\begin{figure}[!htbp]
    \centering
    \includegraphics[width=0.78\linewidth]{experiments/graphs/results/snr_vs_ber.png}
    \caption{SNR vs.\ measured pre-correction BER, with the BPSK reference
    $\tfrac{1}{2}\mathrm{erfc}(\sqrt{\mathrm{SNR}})$ overlaid. The gap between
    the two curves is the effect of the 8D multidimensional-reconciliation
    preprocessor and underlies the empirical SNR$\leftrightarrow$BER mapping
    used in the results chapter.}
    \label{fig:snr_vs_ber_appendix}
\end{figure}


\chapter{Example Experiment Configuration}
\label{chap:appendix-config}

This appendix gives a representative configuration for the external LDPC sweep
runner described in Section~\ref{subsec:meth-orchestration}. The listing is
illustrative and shows the main parameters used to define an SNR sweep.

\begin{lstlisting}[style=appendixcode,language={},caption={Example YAML configuration for an LDPC SNR sweep.},label={lst:yaml-config}]
experiment_name: "snr_sweep_0_to_7_5_1000frames_per_point"
seed: 123

snr:
  start: 0.0001
  end: 7.5
  step_factor: 1.1

frames_per_point: 10000

ldpc:
  binary: "cvqkd-reconciliation/build/ldpc_experiment"
  b_matrices_dir: "cvqkd-reconciliation/data/B_matrices"
  b_matrices: ["256x255_z1.coo", "512x511_z4.coo", "1024x1023_z1.coo"]
  qkd_dimension: 8

output:
  raw_dir: "experiments/data/raw"
\end{lstlisting}


\chapter{Tracked Statistics and Output Format}
\label{chap:appendix-output}

This appendix summarises the structured statistics recorded for both
implementations and gives representative output examples. The LDPC tracking
layer was designed to align closely with the output style of
\texttt{Cascade-cpp}.

\section{Cascade tracked statistics}

\begin{table}[!htbp]
    \centering
    \small
    \setlength{\tabcolsep}{4pt}
    \renewcommand{\arraystretch}{1.05}
    \begin{tabular}{p{0.45\linewidth} p{0.5\linewidth}}
        \hline
        Field & Meaning \\
        \hline
        \texttt{elapsed\_process\_time} & CPU process time used by the reconciliation run. \\
        \texttt{elapsed\_real\_time} & Wall-clock runtime of the reconciliation run. \\
        \texttt{normal\_iterations} & Number of normal Cascade iterations. \\
        \texttt{biconf\_iterations} & Number of BICONF iterations, where applicable. \\
        \texttt{start\_iteration\_messages} & Messages used to begin iterations. \\
        \texttt{start\_iteration\_bits} & Bits used in iteration-start metadata, including shuffle-related information. \\
        \texttt{ask\_parity\_messages} & Number of parity round-trips. \\
        \texttt{ask\_parity\_blocks} & Total number of blocks for which parity was requested. \\
        \texttt{ask\_parity\_bits} & Total number of bits used for parity requests. \\
        \texttt{reply\_parity\_bits} & Total number of bits used in parity replies. \\
        \texttt{reconciliation\_bits} & Total classical bits exchanged during reconciliation. \\
        \texttt{efficiency} & Reconciliation efficiency relative to the Shannon limit. \\
        \texttt{reconciliation\_bits
        \_per\_key\_bit} & Classical communication overhead normalised per key bit. \\
        \texttt{infer\_parity\_blocks} & Number of block parities inferred rather than explicitly requested. \\
        \hline
    \end{tabular}
    \caption{Main statistics recorded by the \texttt{Cascade-cpp} implementation.}
    \label{tab:Cascade-tracking-fields}
\end{table}

\section{LDPC tracked statistics}

\begin{table}[!htbp]
    \centering
    \small
    \setlength{\tabcolsep}{4pt}
    \renewcommand{\arraystretch}{1.05}
    \begin{tabular}{p{0.45\linewidth} p{0.50\linewidth}}
        \hline
        Field & Meaning \\
        \hline
        \texttt{algorithm\_name} & Algorithm identifier, always \texttt{LDPC}. \\
        \texttt{key\_size} & Number of bits processed per frame. \\
        \texttt{reconciliations} & Number of independent frames processed. \\
        \texttt{snr} & Operating signal-to-noise ratio for the run. \\
        \texttt{actual\_bit\_errors} & Mean and deviation of pre-correction bit errors. \\
        \texttt{actual\_bit\_error\_rate} & Mean and deviation of pre-correction BER. \\
        \texttt{remaining\_bit\_errors} & Mean and deviation of post-correction bit errors. \\
        \texttt{remaining\_bit\_error\_rate} & Mean and deviation of post-correction BER. \\
        \texttt{remaining\_frame
        \_error\_rate} & Mean and deviation of post-correction FER. \\
        \texttt{elapsed\_process\_time} & Mean and deviation of CPU process time. \\
        \texttt{elapsed\_real\_time} & Mean and deviation of wall-clock runtime. \\
        \texttt{normal\_iterations} & Mean and deviation of Min-Sum decoder iterations. \\
        \texttt{reconciliation\_bits} & Mean and deviation of syndrome size or disclosed reconciliation bits. \\
        \texttt{reconciliation\_bits
        \_per\_key\_bit} & Mean and deviation of overhead normalised per key bit. \\
        \texttt{efficiency} & Mean and deviation of reconciliation efficiency, calculated compatibly with Cascade. \\
        \hline
    \end{tabular}
    \caption{Main statistics recorded by the LDPC tracking extension.}
    \label{tab:ldpc-tracking-fields}
\end{table}

\clearpage
\section{Example structured output}

\begin{lstlisting}[style=appendixcode,language={},caption={Illustrative NDJSON-style summary record for one LDPC operating point.},label={lst:ndjson-example}]
{
  "algorithm_name": "LDPC",
  "key_size": 512,
  "reconciliations": 10000,
  "snr": 0.123,
  "actual_bit_error_rate_mean": 0.0812,
  "actual_bit_error_rate_std": 0.0031,
  "remaining_bit_error_rate_mean": 0.0004,
  "remaining_bit_error_rate_std": 0.0012,
  "remaining_frame_error_rate_mean": 0.0170,
  "remaining_frame_error_rate_std": 0.1293,
  "elapsed_process_time_mean": 2.1e-5,
  "elapsed_process_time_std": 4.8e-6,
  "elapsed_real_time_mean": 2.4e-5,
  "elapsed_real_time_std": 5.1e-6,
  "normal_iterations_mean": 18.6,
  "normal_iterations_std": 7.4,
  "reconciliation_bits_mean": 256.0,
  "reconciliation_bits_std": 0.0,
  "reconciliation_bits_per_key_bit_mean": 0.5,
  "reconciliation_bits_per_key_bit_std": 0.0,
  "efficiency_mean": 1.084,
  "efficiency_std": 0.019
}
\end{lstlisting}

\printbibliography

\end{document}